\documentclass[10pt]{article}

\usepackage{amssymb,amscd,amsthm, verbatim,amsmath,color,fancyhdr, mathrsfs}
\usepackage{graphicx}
\usepackage{turnstile}
\usepackage[hidelinks]{hyperref}

\def\tequiv{\ensuremath\sdststile{}{}}

\usepackage[a4paper, left=2cm, right=2cm, top=3cm,
bottom=3cm,dvips]{geometry}

\newtheorem{thm}{Theorem}[section]
\newtheorem{prop}[thm]{Proposition}
\newtheorem{lem}[thm]{Lemma}
\newtheorem{rem}[thm]{Remark}
\newtheorem{cor}[thm]{Corollary}
\newtheorem{exr}[thm]{Exercise}
\newtheorem{ex}[thm]{Example}
\newtheorem{defn}[thm]{Definition}

\newcommand{\R}{\mathbb{R}}
\newcommand{\Z}{\mathbb{Z}}
\newcommand{\N}{\mathbb{N}}
\newcommand{\E}{\mathbb{E}}
\newcommand{\V}{\mathbb{V}}
\newcommand{\K}{\mathbb{K}}

\newcommand{\aN}{\mathcal{N}}
\newcommand{\aB}{\mathcal{B}}
\newcommand{\aM}{\mathcal{M}}

\DeclareMathOperator{\cov}{cov}
\DeclareMathOperator{\Var}{Var}
\DeclareMathOperator{\corr}{corr}

\title{Vecteurs Gaussiens}

\author{Victor Schmit}

\date{\today}

\begin{document}
\maketitle

\tableofcontents
\newpage

\section{Noyau de transition}

\subsection{Généralités}

\subsection{Cas des vecteurs gaussiens}

Soient $X$ et $Z$ deux vecteurs aléatoires gaussiens, à valeurs respectivement dans $\R^n$ et $\R^m$.
L'objectif est de donner un sens à la loi conditionnelle $Z \mid X=x$ pour $x \in \R^n$.

\medskip

On suppose, pour simplifier, que le vecteur $(X,Z)$ est centré et admet une matrice de covariance inversible
\[
	\Sigma_{(X,Z)} =
	\begin{pmatrix}
		\Sigma_X    & \Sigma_{XZ} \\
		\Sigma_{ZX} & \Sigma_Z
	\end{pmatrix}.
\]
Comme $\Sigma_{(X,Z)}$ est définie positive, les matrices $\Sigma_X$ et $\Sigma_Z$ sont elles aussi définies positives, donc inversibles. En effet, si $\Sigma_X$ n'était pas inversible, il existerait $u \in \R^n \setminus \{0\}$ tel que
\[
	u^\top \Sigma_X u = 0.
\]
En posant
\[
	y = \binom{u}{0} \in \R^{n+m},
\]
on obtiendrait
\[
	y^\top \Sigma_{(X,Z)} y = u^\top \Sigma_X u = 0,
\]
ce qui contredit la définie-positivité de $\Sigma_{(X,Z)}$.

\medskip

Les vecteurs $X$ et $Z$ admettent donc des densités données par
\[
	f_X(x)
	=
	\frac{1}{(2\pi)^{n/2}\det(\Sigma_X)^{1/2}}
	\exp\left(-\frac12 x^\top \Sigma_X^{-1} x\right),
\]
et
\[
	f_Z(z)
	=
	\frac{1}{(2\pi)^{m/2}\det(\Sigma_Z)^{1/2}}
	\exp\left(-\frac12 z^\top \Sigma_Z^{-1} z\right).
\]

On définit alors la densité conditionnelle de $Z$ sachant $X=x$ par
\[
	f_{Z\mid X=x}(z) = \frac{f_{(X,Z)}(x,z)}{f_X(x)}.
\]
On peut donc définir le noyau de transition $Q : \mathcal B(\R^m)\times \R^n \to \R_+$ par
\[
	Q(A,x)=\int_A f_{Z\mid X=x}(z)\,dz.
\]
Alors
\[
	\mathbb P(Z\in A\mid X)=Q(A,X).
\]

\begin{thm}
	Pour tout $x\in\R^n$, la loi conditionnelle $Z\mid X=x$ est gaussienne, de la forme
	\[
		Z\mid X=x \sim \mathcal N(\mu(x),\Sigma_c),
	\]
	avec
	\[
		\mu(x)=\Sigma_{ZX}\Sigma_X^{-1}x,
		\qquad
		\Sigma_c=\Sigma_Z-\Sigma_{ZX}\Sigma_X^{-1}\Sigma_{XZ}.
	\]
\end{thm}

\begin{proof}
	Comme $(X,Z)$ est un vecteur gaussien, il existe une matrice $A\in \mathcal M_{m,n}(\R)$ telle que
	\[
		Z=AX+\varepsilon,
	\]
	où $\varepsilon$ est un vecteur gaussien indépendant de $X$.

	En prenant la covariance avec $X$, on obtient
	\[
		\Sigma_{ZX}=\cov(Z,X)=A\Sigma_X,
	\]
	d'où
	\[
		A=\Sigma_{ZX}\Sigma_X^{-1}.
	\]

	Par conséquent,
	\[
		Z=\Sigma_{ZX}\Sigma_X^{-1}X+\varepsilon.
	\]
	Conditionnellement à $X=x$, on a donc
	\[
		Z\mid X=x = \Sigma_{ZX}\Sigma_X^{-1}x+\varepsilon.
	\]
	Comme $\varepsilon$ est gaussien indépendant de $X$, la loi conditionnelle de $Z\mid X=x$ est gaussienne, de moyenne
	\[
		\mu(x)=\Sigma_{ZX}\Sigma_X^{-1}x
	\]
	et de matrice de covariance
	\[
		\Sigma_c=\cov(\varepsilon).
	\]

	Il reste à calculer $\Sigma_c$. Comme
	\[
		\varepsilon = Z-AX,
	\]
	on a
	\[
		\Sigma_c
		=
		\cov(Z-AX)
		=
		\Sigma_Z-A\Sigma_{XZ}-\Sigma_{ZX}A^\top + A\Sigma_XA^\top.
	\]
	En remplaçant $A$ par $\Sigma_{ZX}\Sigma_X^{-1}$, on obtient
	\[
		\Sigma_c
		=
		\Sigma_Z-\Sigma_{ZX}\Sigma_X^{-1}\Sigma_{XZ}.
	\]
	Cela conclut la preuve.
\end{proof}

\begin{rem}
	La matrice de covariance conditionnelle $\Sigma_c$ ne dépend pas de $x$. Dans le cadre gaussien, tous les points $x$ apportent donc la même quantité d'information sur la dispersion résiduelle de $Z$.
\end{rem}

\begin{rem}
	La formule
	\[
		\Sigma_c=\Sigma_Z-\Sigma_{ZX}\Sigma_X^{-1}\Sigma_{XZ}
	\]
	montre que l'information apportée par l'observation de $X$ réduit la variance : la covariance conditionnelle est plus petite que la covariance initiale au sens de l'ordre de Loewner.
\end{rem}

\section{Graphical Lasso}

\subsection{Problème}
On s'intéresse au problème suivant : on observe $n$ sources
$(X_p)_{1 \leq p \leq n}$ qui émettent simultanément.
La question est de savoir comment mesurer les dépendances entre ces différentes variables aléatoires.

Pour cela, on adopte un modèle gaussien et on suppose que
\[
	X=(X_1,\dots,X_n)^\top \sim \mathcal N(0,\Sigma),
\]
où $\Sigma$ désigne la matrice de covariance. On note
\[
	\Theta = \Sigma^{-1}
\]
la matrice de précision.

L'objectif est de quantifier la dépendance résiduelle entre deux variables $X_j$ et $X_k$ lorsque toutes les autres variables sont observées.

\medskip

Pour cela, on note
\[
	X_{-(jk)}=(X_p)_{p\neq j,k}.
\]
On considère les résidus obtenus après régression de $X_j$ et $X_k$ sur les autres variables :
\[
	R_j = X_j - \E(X_j\mid X_{-(jk)}),
	\qquad
	R_k = X_k - \E(X_k\mid X_{-(jk)}).
\]
La \emph{corrélation partielle} entre $X_j$ et $X_k$ sachant les autres variables est alors définie par
\[
	\rho_{jk\mid -(jk)} = \corr(R_j,R_k).
\]

\begin{rem}
	Cette définition mesure la dépendance entre les composantes de $X_j$ et $X_k$ qui ne sont pas expliquées par les autres variables.
\end{rem}

\begin{rem}
	En général, la covariance conditionnelle
	\[
		\cov(X_j,X_k\mid X_{-(jk)})
	\]
	est une variable aléatoire. Dans le cadre gaussien, elle ne dépend pas de la valeur de $X_{-(jk)}$, ce qui permet de la relier à un coefficient réel de corrélation partielle.
\end{rem}

\begin{rem}
	Si l'on écrit la régression linéaire
	\[
		X_j = \beta_k X_k + \sum_{r\neq j,k}\beta_r X_r + \varepsilon_j,
	\]
	alors le résidu $\varepsilon_j$ est orthogonal aux variables explicatives. Dans le cadre gaussien, il leur est même indépendant.
\end{rem}

\subsection{Complément de Schur}

Pour poursuivre l'étude sur la corrélation, on va avoir besoin de quelques résultats pour inverser les matrices par blocs.

\begin{defn}[Complément de Schur]
	Soit $M \in \aM_{n+m}(\R)$ une matrice. Alors on peut écrire
	\[
		M =
		\begin{pmatrix}
			A & B \\
			C & D \\
		\end{pmatrix}.
	\]
	Si $D$ est inversible on note $S = A - BD^{-1}C$, le complément de Schur.
\end{defn}

\begin{rem}
	Le complément de Schur peut s'apparenter à une élimination de Gauss pour une matrice de taille $2 \times 2$.
	En effet, on veut résoudre le système
	\[
		M
		\begin{pmatrix}
			x \\
			y
		\end{pmatrix}
		=
		\begin{pmatrix}
			u \\
			v
		\end{pmatrix}.
	\]
	Donc,
	\begin{align}
		A x + B y & = u \\
		C x + D y & = v
	\end{align}
	Si $(1) \to (1) - BD^{-1} (2)$, on a alors
	\begin{align}
		(A - B D^{-1} C) x & = u - B D^{-1}v \\
		C x + D y          & = v
	\end{align}
	D'un point de vue matricielle, cela veut dire que
	\[
		\begin{pmatrix}
			A & B \\
			C & D
		\end{pmatrix}
		\begin{pmatrix}
			I_n      & 0   \\
			-D^{-1}C & I_m
		\end{pmatrix} =
		\begin{pmatrix}
			S & B \\
			0 & D
		\end{pmatrix}.
	\]
	On a donc le résultat suivant.
\end{rem}

\begin{prop}
	La matrice $M$ est inversible si et seulement si $S$ et $D$ sont inversibles.
	Dans ce cas
	\[
		M^{-1} =
		\begin{pmatrix}
			I_n       & 0   \\
			-D^{-1} C & I_m
		\end{pmatrix}
		\begin{pmatrix}
			S^{-1} & 0      \\
			0      & D^{-1}
		\end{pmatrix}
		\begin{pmatrix}
			I_n & -BD^{-1} \\
			0   & I_m
		\end{pmatrix}.
	\]
	Ou encore
	\[
		M^{-1} =
		\begin{pmatrix}
			S^{-1}           & -S^{-1}BD^{-1} \\
			-D^{-1} C S^{-1} & D^{-1}
		\end{pmatrix}.
	\]
\end{prop}

\subsection{Algorithme}

\begin{prop}
	Pour tout $j\neq k$, on a
	\[
		\rho_{jk\mid -(jk)}
		=
		-\frac{\Theta_{jk}}{\sqrt{\Theta_{jj}\Theta_{kk}}}.
	\]
	En particulier,
	\[
		\rho_{jk\mid -(jk)}=0
		\quad \Longleftrightarrow \quad
		\Theta_{jk}=0.
	\]
	Ainsi, dans le modèle gaussien, l'absence d'arête entre les sommets $j$ et $k$ dans le graphe de dépendance conditionnelle équivaut à une entrée nulle dans la matrice de précision.
\end{prop}


\begin{proof}
	Fixons $j\neq k$ et posons
	\[
		A=\{j,k\}, \qquad B=\{1,\dots,n\}\setminus A.
	\]
	On partitionne la matrice de covariance sous la forme
	\[
		\Sigma =
		\begin{pmatrix}
			\Sigma_{AA} & \Sigma_{AB} \\
			\Sigma_{BA} & \Sigma_{BB}
		\end{pmatrix}.
	\]
	Dans le cas gaussien, la loi conditionnelle de $X_A=(X_j,X_k)^\top$ sachant $X_B$ est gaussienne de covariance
	\[
		\Sigma_{A\mid B}
		=
		\Sigma_{AA}-\Sigma_{AB}\Sigma_{BB}^{-1}\Sigma_{BA}.
	\]
	Par la formule du complément de Schur, on a aussi
	\[
		\Sigma_{A\mid B} = (\Theta_{AA})^{-1},
	\]
	où
	\[
		\Theta_{AA}=
		\begin{pmatrix}
			\Theta_{jj} & \Theta_{jk} \\
			\Theta_{jk} & \Theta_{kk}
		\end{pmatrix}.
	\]
	On en déduit
	\[
		(\Theta_{AA})^{-1}
		=
		\frac{1}{\Theta_{jj}\Theta_{kk}-\Theta_{jk}^2}
		\begin{pmatrix}
			\Theta_{kk}  & -\Theta_{jk} \\
			-\Theta_{jk} & \Theta_{jj}
		\end{pmatrix}.
	\]
	Ainsi,
	\[
		\cov(X_j,X_k\mid X_B)
		=
		-\frac{\Theta_{jk}}{\Theta_{jj}\Theta_{kk}-\Theta_{jk}^2},
	\]
	et
	\[
		\Var(X_j\mid X_B)
		=
		\frac{\Theta_{kk}}{\Theta_{jj}\Theta_{kk}-\Theta_{jk}^2},
		\qquad
		\Var(X_k\mid X_B)
		=
		\frac{\Theta_{jj}}{\Theta_{jj}\Theta_{kk}-\Theta_{jk}^2}.
	\]
	Par conséquent,
	\[
		\rho_{jk\mid -(jk)}
		=
		\frac{\cov(X_j,X_k\mid X_B)}
		{\sqrt{\Var(X_j\mid X_B)\Var(X_k\mid X_B)}}
		=
		-\frac{\Theta_{jk}}{\sqrt{\Theta_{jj}\Theta_{kk}}}.
	\]
	Cela prouve la formule.
\end{proof}

Le \emph{Graphical Lasso} consiste alors à estimer une matrice de précision parcimonieuse en résolvant un problème de la forme
\[
	\widehat{\Theta}
	=
	\arg\min_{\Theta \succ 0}
	\left(
	-\log\det(\Theta)
	+
	\operatorname{tr}(\widehat{\Sigma}\Theta)
	+
	\lambda \sum_{j\neq k} |\Theta_{jk}|
	\right),
\]
où $\widehat{\Sigma}$ est la covariance empirique et $\lambda>0$ est un paramètre de régularisation favorisant la parcimonie des coefficients hors diagonale.

\begin{rem}
	La pénalisation en norme $\ell^1$ sur les coefficients hors diagonale favorise l'apparition de zéros dans $\widehat{\Theta}$, et donc l'estimation d'un graphe de dépendance conditionnelle parcimonieux.
\end{rem}



\end{document}
